\begin{table}[!htbp]
\centering
\caption{Pontos de virada das conclusões sob erro hipotético da referência operacional, análise exploratória \textit{post hoc}.}
\label{tab:cap8-vies-referencia}
\small
\begin{tabular}{p{0.25\linewidth}p{0.65\linewidth}}
\toprule
Conclusão & Sensibilidade ao erro da referência \\
\midrule
C1: VPP do LLM inferior a 50\% na vigilância & Igualdade em 40,5\% dos FP reclassificados uniformemente como vínculos, sem retirada de PAR~=~2; primeira virada na grade: 50\%. \\
C2: 13 e 27 perdas & Primeira mudança na grade com 5\% dos PAR~=~2 retirados; com 20\%, perdas esperadas de 10,4 e 21,6. \\
C3: margem de 25\% de O3 frente a O1, meta nominal de 64\% & Sem virada pela mediana nos 100 cenários da vigilância; P97,5 cruza zero em um cenário. \\
C4: F4 exige menos revisões que O1, metas nominais de 25\% e 100\% & Sem virada na grade, nos dois modos. \\
\bottomrule
\end{tabular}
\par\footnotesize Fonte: \path{data/vies_referencia_viradas.csv} e \path{data/vies_referencia.csv}. C1 contínua: $(0{,}5\times25-4)/21$, solução algébrica, não prevalência estimada de erro. C2: $13\times0{,}8$ e $27\times0{,}8$, sob retirada uniforme. Para C3 e C4, as metas são recalculadas após a mudança da referência; mantêm-se as filas, sem retreinamento. Medianas e percentis descrevem 1.000 sorteios por cenário, não intervalos de confiança inferenciais. C1 não é critério de segurança para aceitação autônoma.
\end{table}
